% Preamble (LNCS-compatible):
% \usepackage{booktabs}
% \usepackage{xcolor}
% \usepackage{mdframed}
% \definecolor{verylightpink}{HTML}{FFF5F8}
% \mdfsetup{framemethod=tikz}

%------------------------------------------------------------
% Audit constraints (c1.*)
%------------------------------------------------------------

% !TEX root = ../main.tex

\begin{table}[t]
\small
\setlength{\tabcolsep}{3pt}
\renewcommand{\arraystretch}{1.05}
\caption{Audit column constraints (c1.*).}
\label{tab:constraints-audit}
\begin{tabular}{@{}p{0.98\linewidth}@{}}
\toprule

c1.1 Audit is zero-padded. For all $X\in H$: $\mathsf{Audit}(X)\cdot \mathsf{Z}_{\mathsf{head}}(X)=0$.
\begin{mdframed}[linewidth=0.6pt,backgroundcolor=verylightpink,innerleftmargin=6pt,innerrightmargin=6pt,innertopmargin=4pt,innerbottommargin=4pt]
\[
V_{c1.1}(X)\ :=\ \mathsf{Audit}(X)\,\mathsf{Z}_{\mathsf{head}}(X)
\]
\end{mdframed}
\\[-2mm]
\midrule

c1.2 Audit is binary. For all $X\in H$: $\mathsf{Audit}(X)\bigl(\mathsf{Audit}(X)-1\bigr)=0$.
\begin{mdframed}[linewidth=0.6pt,backgroundcolor=verylightpink,innerleftmargin=6pt,innerrightmargin=6pt,innertopmargin=4pt,innerbottommargin=4pt]
\[
V_{c1.2}(X)\ :=\ \mathsf{Audit}(X)\bigl(\mathsf{Audit}(X)-1\bigr)
\]
\end{mdframed}
\\[-2mm]
\midrule

c1.3 Audit is block-consistent: each ballot block is all-0 or all-1. For all $X\in H$:
$\bigl(\mathsf{Sel}_{\mathsf{blk}}(X)\mathsf{BlkSum}_{\mathsf{Audit}}(X)\bigr)\bigl(\mathsf{Sel}_{\mathsf{blk}}(X)\mathsf{BlkSum}_{\mathsf{Audit}}(X)-\mathcal{C}\bigr)=0$,
where $\mathsf{BlkSum}_{\mathsf{Audit}}(X)=\sum_{k=0}^{\mathcal{C}-1}\mathsf{Audit}(\omega^k X)$.
\begin{mdframed}[linewidth=0.6pt,backgroundcolor=verylightpink,innerleftmargin=6pt,innerrightmargin=6pt,innertopmargin=4pt,innerbottommargin=4pt]
\[
V_{c1.3}(X)\ :=\ \bigl(\mathsf{Sel}_{\mathsf{blk}}(X)\mathsf{BlkSum}_{\mathsf{Audit}}(X)\bigr)\bigl(\mathsf{Sel}_{\mathsf{blk}}(X)\mathsf{BlkSum}_{\mathsf{Audit}}(X)-\mathcal{C}\bigr)
\]
\end{mdframed}
\\[-2mm]
\midrule

c1.4a Audit accumulator step (suffix sum). For all $X\in H\setminus\{\omega^{n-1}\}$:
$\mathsf{Acc}_{\mathsf{Audit}}(X)-\mathsf{Audit}(X)+\mathsf{Acc}_{\mathsf{Audit}}(\omega X)=0$.
\begin{mdframed}[linewidth=0.6pt,backgroundcolor=verylightpink,innerleftmargin=6pt,innerrightmargin=6pt,innertopmargin=4pt,innerbottommargin=4pt]
\[
V_{c1.4a}(X)\ :=\ \bigl(\mathsf{Acc}_{\mathsf{Audit}}(X)-\mathsf{Audit}(X)+\mathsf{Acc}_{\mathsf{Audit}}(\omega X)\bigr)\cdot(X-\omega^{n-1})
\]
\end{mdframed}
\\[-2mm]
\midrule

c1.4b Audit accumulator base (last row). For $X=\omega^{n-1}$:
$\mathsf{Acc}_{\mathsf{Audit}}(X)-\mathsf{Audit}(X)=0$.
\begin{mdframed}[linewidth=0.6pt,backgroundcolor=verylightpink,innerleftmargin=6pt,innerrightmargin=6pt,innertopmargin=4pt,innerbottommargin=4pt]
\[
V_{c1.4b}(X)\ :=\ \bigl(\mathsf{Acc}_{\mathsf{Audit}}(X)-\mathsf{Audit}(X)\bigr)\cdot \mathsf{Mask}_{n-1}(X)
\]
\end{mdframed}
\\[-2mm]
\midrule

c1.4c Audit sum matches public count. For $X=\omega^0$:
$\mathsf{Acc}_{\mathsf{Audit}}(X)-\mathsf{Sum}_{\mathsf{Audit}\times\mathcal{C}}=0$.
\begin{mdframed}[linewidth=0.6pt,backgroundcolor=verylightpink,innerleftmargin=6pt,innerrightmargin=6pt,innertopmargin=4pt,innerbottommargin=4pt]
\[
V_{c1.4c}(X)\ :=\ \bigl(\mathsf{Acc}_{\mathsf{Audit}}(X)-\mathsf{Sum}_{\mathsf{Audit}\times\mathcal{C}}\bigr)\cdot \mathsf{Mask}_{0}(X)
\]
\end{mdframed}
\\[-2mm]

\bottomrule
\end{tabular}
\end{table}

%------------------------------------------------------------
% Mark constraints (c2.*)
%------------------------------------------------------------
\begin{table}[t]
\small
\setlength{\tabcolsep}{3pt}
\renewcommand{\arraystretch}{1.05}
\caption{Mark column constraints (c2.*).}
\label{tab:constraints-mark}
\begin{tabular}{@{}p{0.98\linewidth}@{}}
\toprule

c2.1 Mark is zero-padded. For all $X\in H$: $\mathsf{Mark}(X)\cdot \mathsf{Z}_{\mathsf{head}}(X)=0$.
\begin{mdframed}[linewidth=0.6pt,backgroundcolor=verylightpink,innerleftmargin=6pt,innerrightmargin=6pt,innertopmargin=4pt,innerbottommargin=4pt]
\[
V_{c2.1}(X)\ :=\ \mathsf{Mark}(X)\,\mathsf{Z}_{\mathsf{head}}(X)
\]
\end{mdframed}
\\[-2mm]
\midrule

c2.2 Mark is binary. For all $X\in H$: $\mathsf{Mark}(X)\bigl(\mathsf{Mark}(X)-1\bigr)=0$.
\begin{mdframed}[linewidth=0.6pt,backgroundcolor=verylightpink,innerleftmargin=6pt,innerrightmargin=6pt,innertopmargin=4pt,innerbottommargin=4pt]
\[
V_{c2.2}(X)\ :=\ \mathsf{Mark}(X)\bigl(\mathsf{Mark}(X)-1\bigr)
\]
\end{mdframed}
\\[-2mm]
\midrule

c2.3 Mark has at most one ‘1’ per ballot block. For all $X\in H$:
$\bigl(\mathsf{Sel}_{\mathsf{blk}}(X)\mathsf{BlkSum}_{\mathsf{Mark}}(X)\bigr)\bigl(\mathsf{Sel}_{\mathsf{blk}}(X)\mathsf{BlkSum}_{\mathsf{Mark}}(X)-1\bigr)=0$,
where $\mathsf{BlkSum}_{\mathsf{Mark}}(X)=\sum_{k=0}^{\mathcal{C}-1}\mathsf{Mark}(\omega^k X)$.
\begin{mdframed}[linewidth=0.6pt,backgroundcolor=verylightpink,innerleftmargin=6pt,innerrightmargin=6pt,innertopmargin=4pt,innerbottommargin=4pt]
\[
V_{c2.3}(X)\ :=\ \bigl(\mathsf{Sel}_{\mathsf{blk}}(X)\mathsf{BlkSum}_{\mathsf{Mark}}(X)\bigr)\bigl(\mathsf{Sel}_{\mathsf{blk}}(X)\mathsf{BlkSum}_{\mathsf{Mark}}(X)-1\bigr)
\]
\end{mdframed}
\\[-2mm]
\midrule

c2.4a Mark accumulator step (suffix sum). For all $X\in H\setminus\{\omega^{n-1}\}$:
$\mathsf{Acc}_{\mathsf{Mark}}(X)-\mathsf{Mark}(X)+\mathsf{Acc}_{\mathsf{Mark}}(\omega X)=0$.
\begin{mdframed}[linewidth=0.6pt,backgroundcolor=verylightpink,innerleftmargin=6pt,innerrightmargin=6pt,innertopmargin=4pt,innerbottommargin=4pt]
\[
V_{c2.4a}(X)\ :=\ \bigl(\mathsf{Acc}_{\mathsf{Mark}}(X)-\mathsf{Mark}(X)+\mathsf{Acc}_{\mathsf{Mark}}(\omega X)\bigr)\cdot(X-\omega^{n-1})
\]
\end{mdframed}
\\[-2mm]
\midrule

c2.4b Mark accumulator base (last row). For $X=\omega^{n-1}$:
$\mathsf{Acc}_{\mathsf{Mark}}(X)-\mathsf{Mark}(X)=0$.
\begin{mdframed}[linewidth=0.6pt,backgroundcolor=verylightpink,innerleftmargin=6pt,innerrightmargin=6pt,innertopmargin=4pt,innerbottommargin=4pt]
\[
V_{c2.4b}(X)\ :=\ \bigl(\mathsf{Acc}_{\mathsf{Mark}}(X)-\mathsf{Mark}(X)\bigr)\cdot \mathsf{Mask}_{n-1}(X)
\]
\end{mdframed}
\\[-2mm]
\midrule

c2.4c Mark sum matches public count. For $X=\omega^0$:
$\mathsf{Acc}_{\mathsf{Mark}}(X)-\mathsf{Sum}_{\mathsf{Mark}}=0$.
\begin{mdframed}[linewidth=0.6pt,backgroundcolor=verylightpink,innerleftmargin=6pt,innerrightmargin=6pt,innertopmargin=4pt,innerbottommargin=4pt]
\[
V_{c2.4c}(X)\ :=\ \bigl(\mathsf{Acc}_{\mathsf{Mark}}(X)-\mathsf{Sum}_{\mathsf{Mark}}\bigr)\cdot \mathsf{Mask}_{0}(X)
\]
\end{mdframed}
\\[-2mm]

\bottomrule
\end{tabular}
\end{table}

%------------------------------------------------------------
% Cross-column constraint (c3)
%------------------------------------------------------------
\begin{table}[t]
\small
\setlength{\tabcolsep}{3pt}
\renewcommand{\arraystretch}{1.05}
\caption{Cross-column constraint (c3).}
\label{tab:constraints-cross}
\begin{tabular}{@{}p{0.98\linewidth}@{}}
\toprule

c3 No ballot is both audited and voted. For all $X\in H$: $\mathsf{Audit}(X)\cdot \mathsf{Mark}(X)=0$.
\begin{mdframed}[linewidth=0.6pt,backgroundcolor=verylightpink,innerleftmargin=6pt,innerrightmargin=6pt,innertopmargin=4pt,innerbottommargin=4pt]
\[
V_{c3}(X)\ :=\ \mathsf{Audit}(X)\,\mathsf{Mark}(X)
\]
\end{mdframed}
\\[-2mm]

\bottomrule
\end{tabular}
\end{table}

%------------------------------------------------------------
% Tally constraints (c4.*)
%------------------------------------------------------------
\begin{table}[t]
\small
\setlength{\tabcolsep}{3pt}
\renewcommand{\arraystretch}{1.05}
\caption{Tally constraints (c4.*).}
\label{tab:constraints-tally}
\begin{tabular}{@{}p{0.98\linewidth}@{}}
\toprule

c4.1 Strided vote-accumulator step (stride $\mathcal{C}$). For all
$X\in H\setminus\{\omega^{n-1},\dots,\omega^{n-\mathcal{C}}\}$:
$\mathsf{Acc}_{\mathsf{Vote}}(X)-\mathsf{Mark}(X)+\mathsf{Acc}_{\mathsf{Vote}}(\omega^{\mathcal{C}}X)=0$.
\begin{mdframed}[linewidth=0.6pt,backgroundcolor=verylightpink,innerleftmargin=6pt,innerrightmargin=6pt,innertopmargin=4pt,innerbottommargin=4pt]
\[
V_{c4.1}(X)\ :=\ \bigl(\mathsf{Acc}_{\mathsf{Vote}}(X)-\mathsf{Mark}(X)+\mathsf{Acc}_{\mathsf{Vote}}(\omega^{\mathcal{C}}X)\bigr)\cdot \mathsf{Mask}_{\mathsf{tail}}(X)
\]
\end{mdframed}
\\[-2mm]
\midrule

c4.2 Strided vote-accumulator tail base. For all
$X\in\{\omega^{n-1},\dots,\omega^{n-\mathcal{C}}\}$:
$\mathsf{Acc}_{\mathsf{Vote}}(X)-\mathsf{Mark}(X)=0$.
\begin{mdframed}[linewidth=0.6pt,backgroundcolor=verylightpink,innerleftmargin=6pt,innerrightmargin=6pt,innertopmargin=4pt,innerbottommargin=4pt]
\[
V_{c4.2}(X)\ :=\ \bigl(\mathsf{Acc}_{\mathsf{Vote}}(X)-\mathsf{Mark}(X)\bigr)\cdot \frac{\mathsf{Z}_H(X)}{\mathsf{Mask}_{\mathsf{tail}}(X)}
\]
\end{mdframed}
\\[-2mm]
\midrule

c4.3 Tally readout. For each $c\in\{0,\dots,\mathcal{C}-1\}$, for $X=\omega^{c}$:
$\mathsf{Acc}_{\mathsf{Vote}}(X)-\mathsf{Tally}[c]=0$.
\begin{mdframed}[linewidth=0.6pt,backgroundcolor=verylightpink,innerleftmargin=6pt,innerrightmargin=6pt,innertopmargin=4pt,innerbottommargin=4pt]
\[
V_{c4.3,c}(X)\ :=\ \bigl(\mathsf{Acc}_{\mathsf{Vote}}(X)-\mathsf{Tally}[c]\bigr)\cdot \frac{\mathsf{Z}_H(X)}{X-\omega^{c}}
\]
\end{mdframed}
\\[-2mm]

\bottomrule
\end{tabular}
\end{table}


\section{Voter Checks}

\subsection{Print Audit}

%------------------------------------------------------------
% Notation for voter checks (c5.*): print-audit, receipt, dispute
%------------------------------------------------------------

\paragraph{Indices and ballots.}
We index rows by $i\in\{0,\dots,n-1\}$ and identify the $i$-th domain point with $X=\omega^i$.
A ballot has ID $b\in\{0,\dots,\mathcal{B}-1\}$ and occupies the \emph{ballot block}
\[
I_b=\{\,b\mathcal{C},\,b\mathcal{C}+1,\,\dots,\,b\mathcal{C}+(\mathcal{C}-1)\,\}\subseteq\{0,\dots,n-1\}.
\]

\paragraph{Block selector for ballot $b$.}
Let $\mathsf{Sel}_b(X)$ be the (publicly checkable) selector polynomial for ballot $b$ defined pointwise by
\[
\mathsf{Sel}_b(\omega^i)=
\begin{cases}
1,& i\in I_b,\\
0,& \text{otherwise},
\end{cases}
\qquad i\in\{0,\dots,n-1\}.
\]
(Equivalently, $\mathsf{Sel}_b$ is the interpolation of this $0/1$ array over $H$.)

%------------------------------------------------------------
% c5.1 Print-audit check (open the whole block)
%------------------------------------------------------------
\paragraph{c5.1 Print-audit check.}
The voter supplies ballot ID $b$. The EA returns the $\mathcal{C}$ evaluations
\[
\bigl\{\mathsf{Audit}(\omega^{i})\bigr\}_{i\in I_b},
\qquad
\bigl\{\mathsf{Mark}(\omega^{i})\bigr\}_{i\in I_b},
\qquad
\bigl\{\mathsf{Code}(\omega^{i})\bigr\}_{i\in I_b},
\]
together with a batched opening proof for these evaluations.


\subsection{Receipt Check}

%------------------------------------------------------------
% c5.2 Receipt check (shuffle-based)
%------------------------------------------------------------
\paragraph{c5.2 Receipt check (shuffle-based).}

\emph{Ballot-ID column.}
Let $\mathsf{BallotID}(X)$ be a public column encoding ballot IDs per position, defined by
\[
\mathsf{BallotID}(\omega^{b\mathcal{C}+k})=b \quad\text{for } b\in\{0,\dots,\mathcal{B}-1\},\;k\in\{0,\dots,\mathcal{C}-1\},
\]
and $\mathsf{BallotID}(\omega^i)=0$ on padding indices $i\ge \mathcal{P}$.

\emph{Secret permutation.}
Let $\pi$ be a secret permutation of $\{0,\dots,n-1\}$. For any column $\mathsf{F}$, define its permuted form
\[
\mathsf{F}^{\pi}(\omega^i)=\mathsf{F}\bigl(\omega^{\pi(i)}\bigr).
\]
The EA commits once (pre-election or once-per-election) to
\[
\mathsf{BallotID}^{\pi}(X),\quad \mathsf{Code}^{\pi}(X),\quad \mathsf{Mark}^{\pi}(X),
\]
and proves (once) that each is a permutation of its source column under the \emph{same} $\pi$.

\emph{Per-receipt query.}
Given ballot ID $b$, the EA finds (privately) an index $t$ such that
\[
\mathsf{BallotID}^{\pi}(\omega^t)=b
\quad\text{and}\quad
\mathsf{Mark}^{\pi}(\omega^t)=1,
\]
and returns the claimed receipt code
\[
c := \mathsf{Code}^{\pi}(\omega^t),
\]
along with a batched opening proof at the single point $\omega^t$ for
$\mathsf{BallotID}^{\pi},\mathsf{Mark}^{\pi},\mathsf{Code}^{\pi}$.
(Revealing $t$ is safe because $\pi$ is secret, so $t$ carries no candidate-order meaning.)

%------------------------------------------------------------
% c5.3 Dispute resolution (selector + equality flags)
%------------------------------------------------------------

\subsection{Dispute Resolution}

\paragraph{c5.3 Dispute resolution (selector + equality flags).}
The voter supplies ballot ID $b$ and disputed code $c\in\mathbb{F}_q$.
The EA proves that $c$ does not appear on that ballot block without revealing any codes.

\emph{Equality flag polynomial.}
The EA introduces a binary flag polynomial $\mathsf{Eq}(X)\in\{0,1\}$ over $H$ intended to satisfy
\[
\mathsf{Eq}(\omega^i)=
\begin{cases}
1,& \mathsf{Code}(\omega^i)=c,\\
0,& \text{otherwise}.
\end{cases}
\]

\emph{Inverse witness for correctness.}
The EA also introduces an auxiliary polynomial $\mathsf{Inv}(X)$ and enforces, for all $X\in H$,
\[
(\mathsf{Code}(X)-c)\cdot \mathsf{Inv}(X) - (1-\mathsf{Eq}(X)) = 0,
\]
together with the booleanity constraint
\[
\mathsf{Eq}(X)\bigl(\mathsf{Eq}(X)-1\bigr)=0.
\]
Intuition: if $\mathsf{Code}(X)\neq c$ then $\mathsf{Eq}(X)=0$ and $\mathsf{Inv}(X)=(\mathsf{Code}(X)-c)^{-1}$ exists;
if $\mathsf{Code}(X)=c$ then $\mathsf{Eq}(X)=1$ and the equation forces $0=0$ without requiring an inverse.

\emph{Non-membership on ballot $b$.}
Finally, the EA enforces that no equality occurs inside the ballot block:
\[
\sum_{i=0}^{n-1}\mathsf{Sel}_b(\omega^i)\cdot \mathsf{Eq}(\omega^i)=0.
\]
(Equivalently, define $\mathsf{Hit}(X)=\mathsf{Sel}_b(X)\cdot \mathsf{Eq}(X)$ and prove its sum is $0$.)